\documentclass[11pt, a4paper]{article}

\usepackage[a4paper, top=2.5cm, bottom=2.5cm, left=2cm, right=2cm]{geometry}
\usepackage{amsmath}
\usepackage{amssymb}
\usepackage[colorlinks=true, linkcolor=blue, urlcolor=blue, citecolor=blue]{hyperref}

\title{Theory of Everything - Volume 17 \\ \Large Fluid Dynamics \& Turbulence}
\author{Omega Protocol}
\date{\today}

\begin{document}

\maketitle

\section*{Layman's Summary}
How does water flow? How does air move around a jet wing? Volume 17 connects the everyday movement of fluids to the deepest quantum rules of the universe. We prove that the chaotic swirling of water and the mathematical equations used by engineers are actually just large-scale, "zoomed-out" versions of the universe's foundational informational flow.

\section{The Hydrodynamic Limit (Axiom 20)}
The Navier-Stokes equations have been used for centuries to design everything from submarines to weather forecasts. We show that these equations aren't just useful approximations. If you take the incredibly complex, microscopic quantum informational flow of the universe and "zoom out" (the long-wavelength limit), the Navier-Stokes equations perfectly and naturally emerge.

\section{Viscosity Bound (Theorem)}
Why is honey thicker than water? Fluid viscosity is a measure of how much a liquid resists flowing. Our theory proves that there is a fundamental limit to how "runny" a fluid can be. The ratio of a fluid's thickness to its informational content (entropy) can never drop below a specific number (1/4$\pi$). This connects the physics of a dripping faucet directly to the physics of black holes.

\section{Turbulence Onset (Theorem)}
Turbulence—like the bumpy air during a flight or the frothing white water of a river—is famously one of the hardest problems in physics. We prove that turbulence happens when the smooth "modular flow lines" of quantum information break apart in a fractal pattern. The chaos of a whirlpool is the macroscopic shadow of quantum informational breaking.

\end{document}
